\documentclass[a4paper,11pt]{article}

\usepackage[T1]{fontenc}
\usepackage[utf8]{inputenc}
\usepackage{charter}
\usepackage{latexsym}
\usepackage[left=0.57in,right=0.57in,top=0.42in,bottom=0.4in]{geometry}
\usepackage{titlesec}
\usepackage{marvosym}
\usepackage[usenames,dvipsnames]{color}
\usepackage{enumitem}
\usepackage{fancyhdr}
\usepackage{tabularx}
\usepackage{array}
\usepackage{amssymb}
\usepackage[hidelinks]{hyperref}

\hypersetup{
  pdftitle={Juan Manuel Ruiz Muñoz – CV},
  pdfauthor={Juan Manuel Ruiz Muñoz}
}

\pagestyle{fancy}
\fancyhf{}
\renewcommand{\headrulewidth}{0pt}
\renewcommand{\footrulewidth}{0pt}

\urlstyle{same}
\raggedbottom
\setlength{\tabcolsep}{0in}
\setlength{\parindent}{0pt}

\titleformat{\section}{\vspace{-4pt}\scshape\raggedright\large\bfseries}{}{0em}{}[\color{black}\titlerule \vspace{-3pt}]

\newcommand{\resumeItem}[1]{\item\small{#1 \vspace{-2pt}}}
\newcommand{\resumeSubheading}[4]{
  \vspace{-1pt}\item
    \begin{tabular*}{\textwidth}[t]{l@{\extracolsep{\fill}}r}
      \textbf{#1} & #2 \\
      \textit{\small#3} & \textit{\small #4} \\
    \end{tabular*}\vspace{-5pt}
}
\newcommand{\resumeProjectHeading}[2]{
    \item
    \begin{tabular*}{\textwidth}{l@{\extracolsep{\fill}}r}
      \small#1 & \small#2 \\
    \end{tabular*}\vspace{-5pt}
}
\renewcommand\labelitemii{$\vcenter{\hbox{\tiny$\bullet$}}$}
\newcommand{\resumeSubHeadingListStart}{\begin{itemize}[leftmargin=0in, label={}]}
\newcommand{\resumeSubHeadingListEnd}{\end{itemize}}
\newcommand{\resumeItemListStart}{\begin{itemize}[leftmargin=0.18in, itemsep=0pt, topsep=1pt, parsep=0pt]}
\newcommand{\resumeItemListEnd}{\end{itemize}\vspace{-5pt}}

\begin{document}
\frenchspacing

%----------HEADING----------
\begin{center}
  {\LARGE \textbf{JUAN MANUEL RUIZ MUÑOZ}} \\ \vspace{1pt}
  {\large ML Engineer · Computer Vision \& LLMs} \\ \vspace{2pt}
  \small Alicante, España $\diamond$ +34 678 23 12 12 $\diamond$
  \href{mailto:juanmanuelruiz383@gmail.com}{juanmanuelruiz383@gmail.com} $\diamond$
  \href{https://github.com/juanmanuelruizm}{GitHub} $\diamond$
  \href{https://www.linkedin.com/in/juanmanuelruizmunoz/}{LinkedIn} $\diamond$
  \href{https://juanmanuel.petrer.eu}{juanmanuel.petrer.eu} \\ \vspace{1pt}
  \small\textit{Disponibilidad para trabajo presencial, híbrido o remoto}
\end{center}

%----------PERFIL----------
\section{Perfil}
\small{ML Engineer con formación en Física y Matemáticas y más de dos años llevando modelos de Computer Vision y LLMs de la investigación a producción. Construyo sistemas de principio a fin (datos, entrenamiento, evaluación, despliegue) para inspección industrial, incluido un OCR multimodal ajustado con LoRA con un 98,56\,\% de exact match en producción. Rigor de formación: entiendo por qué funciona un modelo antes de desplegarlo. Publicación en DCAI'25.}

%----------EXPERIENCIA----------
\section{Experiencia profesional}
\resumeSubHeadingListStart

  \resumeSubheading
    {Vraia Corp.}{06/2025 -- Actualidad}
    {ML Engineer · Computer Vision \& LLMs}{Elche, Alicante}
    \resumeItemListStart
      \resumeItem{Desarrollo de principio a fin de un \textbf{OCR multimodal} para leer flancos de neumático: fine-tuning de \textbf{GLM-OCR} con \textbf{LoRA} (LLaMA-Factory) en un servidor 2$\times$RTX 3090, etapa de corrección en post-procesado y serving con \textbf{vLLM}: \textbf{98,56\,\% de exact match en producción}.}
      \resumeItem{Diseño y entrenamiento de modelos de \textbf{detección y segmentación} (CNNs, Vision Transformers, YOLO, RF-DETR) para inspección industrial de neumáticos y \textbf{detección de defectos}: \textbf{mAP@50:95 de 0,83} en segmentación multiclase sobre datasets propios de miles de imágenes.}
      \resumeItem{Diagnóstico y mitigación del colapso de máscara en piezas finas anulares en un modelo de segmentación \textbf{RF-DETR} de 7 clases (ventanas y persianas), reponderando la pérdida de máscara.}
      \resumeItem{Pipeline de \textbf{generación sintética de defectos} con modelos generativos para ampliar clases escasas y mejorar la robustez de la detección.}
      \resumeItem{\textbf{MLOps}: pipelines ETL y de aumento de datos visuales; despliegue y monitorización con \textbf{Docker, MLflow} y TensorBoard; protocolos de evaluación (precision, recall, F1, mAP, IoU); especificación algorítmica en \textbf{C++} para su integración por el equipo del cliente.}
    \resumeItemListEnd

  \resumeSubheading
    {AIR Institute \& Grupo de Investigación BISITE (Universidad de Salamanca)}{06/2024 -- 06/2025}
    {Data Scientist · AI Researcher}{Salamanca}
    \resumeItemListStart
      \resumeItem{Codiseño de una arquitectura \textbf{LLM + RAG} que construye de forma incremental una ontología verificada para detectar afirmaciones falsas en obras de ficción (precisión 0,98; F1 0,86); \textbf{publicada en DCAI'25}.}
      \resumeItem{Pipelines \textbf{RAG} y de recuperación semántica (ChromaDB) y sistemas \textbf{multiagente} (Ollama, DeepSeek, GPT-4o-mini), con foco en reducir coste y latencia.}
      \resumeItem{Proyectos de Computer Vision (detección, clasificación, segmentación) en todo el ciclo de vida ML, en contacto directo con clientes para captura de requisitos y adaptación del estado del arte a casos de uso reales.}
    \resumeItemListEnd

  \resumeSubheading
    {Universidad Miguel Hernández -- Grupo SCOLAb}{01/2024 -- 07/2024}
    {Investigador en prácticas}{Elche, Alicante}
    \resumeItemListStart
      \resumeItem{Métodos lagrangianos para la identificación automática de bloqueos atmosféricos (ERA5, HYSPLIT, DBSCAN); base del TFG y de una publicación en preparación.}
    \resumeItemListEnd

\resumeSubHeadingListEnd

%----------PROYECTO----------
\section{Proyecto destacado}
\resumeSubHeadingListStart
  \resumeProjectHeading
    {\textbf{vcalib} -- filtros de calibración para RF-DETR bajo cambios de iluminación}{\href{https://github.com/juanmanuelruizm/vcalib}{código abierto}}
    \resumeItemListStart
      \resumeItem{Filtro diferenciable de pocos miles de parámetros que corrige el cambio de iluminación \textit{antes} de un detector congelado, sin reentrenar ni etiquetas. En 6 escenas reales no vistas en el ajuste, el \textbf{AP@[.5:.95] pasa de 0,377 a 0,580} con el filtro (0,544 en la condición de referencia), medido contra pseudo-ground-truth independiente de SAM\,3. Se ajusta en $\sim$3,5 min en la CPU de un portátil.}
    \resumeItemListEnd
\resumeSubHeadingListEnd

%----------FORMACION----------
\section{Formación}
\resumeSubHeadingListStart
  \resumeProjectHeading
    {\textbf{Grado en Física} -- Universidad de Alicante}{2020 -- 2024}
    \resumeItemListStart
      \resumeItem{+120 ECTS del Grado en Matemáticas cursados en paralelo (programa de estudios simultáneos de la UA). Beca SICUE en la Universidad de Salamanca (2022 -- 2023).}
    \resumeItemListEnd
\resumeSubHeadingListEnd

%----------PUBLICACIONES----------
\section{Publicaciones}
\resumeSubHeadingListStart
  \vspace{-2pt}
  \item
    \begin{itemize}[leftmargin=0.18in]
      \resumeItem{\textbf{DCAI'25}: \textit{AI-Powered Ontology-Based Architecture for Misinformation Detection in Fiction Works.} \href{https://juanmanuel.petrer.eu/DCAI.pdf}{[PDF]}}
    \end{itemize}
\resumeSubHeadingListEnd
\vspace{-8pt}

%----------COMPETENCIAS----------
\section{Competencias}
\small
\begin{tabularx}{\textwidth}{@{}p{1.21in}>{\raggedright\arraybackslash}X@{}}
  \textbf{LLMs \& NLP} & Fine-tuning LoRA (LLaMA-Factory), Transformers, RAG, vLLM, Ollama, BD vectoriales (ChromaDB, FAISS), multiagente \\
  \textbf{Computer Vision} & PyTorch, OpenCV, CNNs, ViTs, YOLO, RF-DETR, SAM, modelos generativos, detección, segmentación \\
  \textbf{MLOps y código} & Python, C++ (especificación), R, Docker, MLflow, TensorBoard, Git, AWS, Linux, uv; estadística \\
  \textbf{Idiomas} & Español (nativo) · Inglés (uso profesional) \\
\end{tabularx}

\end{document}
